\chapter{Least Action, Local Equations, and Symmetry}

The laws of physics have many subjects but a remarkably common form. Newtonian
particles, electromagnetic fields, gravitation, and even the fields of modern
physics can all be described by an action principle. Thermodynamics looks
different because it is statistical, but its microscopic degrees of freedom
still obey dynamical laws of this kind.

The phrase laws of physics is being used specifically here. Mathematical
identities such as \(1+1=2\), and the internal rules of probability theory, are
not dynamical laws of a physical system. The fact that probability describes
nature may itself be a physical fact, but the laws at issue now are those that
govern motion and fields.

This is a stronger claim than saying that least action is a clever alternative
to \(F=ma\). The action is the organizing language; Newton's equations are one
local expression of it. Our first task is therefore to understand how a global
statement about an entire history can become a differential equation at each
instant of time. Two elementary facts from calculus will do nearly all of the
work.

\section{Two Small Theorems}

\subsection{Integration by parts}

Begin with the fundamental theorem of calculus,
\begin{equation}
 \int_{t_1}^{t_2}\dot F(t)\,dt=F(t_2)-F(t_1).
 \label{eq:l03-fundamental-theorem}
\end{equation}
There is a useful elementary picture behind this formula. Divide the interval
into short steps. Each contribution to the integral is approximately a
difference,
\[
 F(t_{k+1})-F(t_k).
\]
On adding the differences, every intermediate value cancels:
\[
 \bigl(F_2-F_1\bigr)+\bigl(F_3-F_2\bigr)+\cdots
   +\bigl(F_N-F_{N-1}\bigr)=F_N-F_1.
\]
Only the two endpoints survive.

Now take \(F=fg\). The product rule gives
\begin{equation}
 \frac{d}{dt}(fg)=\dot f\,g+f\,\dot g.
\end{equation}
Substitution into Eq.~\eqref{eq:l03-fundamental-theorem} yields
\begin{equation}
 \int_{t_1}^{t_2}\dot f\,g\,dt
 =
 [fg]_{t_1}^{t_2}
 -
 \int_{t_1}^{t_2}f\,\dot g\,dt.
 \label{eq:l03-integration-by-parts}
\end{equation}
Thus a derivative may be moved from one factor to the other, at the cost of a
minus sign and a boundary term.

\begin{figure}[H]
\centering
\includegraphics[width=0.82\textwidth]{lecture_03_figure_05.png}
\caption{Integration by parts on the board at 00:10:20. The upper line retains
the endpoint term; the lower line is being reduced to the endpoint-vanishing
case used in the variation.}
\lectureframe{00:10:20}
\label{fig:l03-integration-by-parts}
\end{figure}

\begin{classroomqa}
\audiencequestion{Why must the integral contain \(dt\)?}
\lecturerresponse{An integral is always taken with respect to a variable. Here
time is the independent variable, so the complete expression is
\(\int \dot f(t)\,dt\). Writing only \(\int\dot f(t)\) leaves the integration
measure unstated. Small notational questions are worth stopping for because
they force every ingredient of the calculation into view.}
\end{classroomqa}

We shall use a special case repeatedly. If \(f\) vanishes at \(t_1\) and
\(t_2\), then \([fg]_{t_1}^{t_2}=0\), and
\begin{equation}
 \int_{t_1}^{t_2}\dot f\,g\,dt
 =
 -\int_{t_1}^{t_2}f\,\dot g\,dt.
 \label{eq:l03-integration-by-parts-fixed}
\end{equation}
It is enough that the product \(fg\) have the same value at both endpoints,
but fixed-endpoint variations will make \(f\) vanish separately at each end.

\subsection{An arbitrary test function}

The second fact is almost obvious, but it is the step that turns an integral
condition into a local law.

\begin{proposition}
If
\begin{equation}
 \int_{t_1}^{t_2}a(t)f(t)\,dt=0
 \label{eq:l03-fundamental-lemma-hypothesis}
\end{equation}
for every sufficiently smooth function \(f(t)\) that vanishes at the
endpoints, then \(a(t)=0\) throughout the interval.
\end{proposition}

\begin{proof}
Choose \(f\) to be a narrow smooth blip around an arbitrary point \(t_*\) and
zero away from that neighborhood. If \(a(t_*)\neq0\), choose the sign of the
blip to agree with the sign of \(a\); for a sufficiently narrow neighborhood
the integral cannot vanish. Hence \(a(t_*)=0\). Since the point was arbitrary,
\(a\) vanishes everywhere.
\end{proof}

\begin{figure}[H]
\centering
\includegraphics[width=0.76\textwidth]{lecture_03_figure_06.png}
\caption{The arbitrary-function argument at 00:13:30. A narrow test function
isolates the value of \(a(t)\) near any chosen point of the interval.}
\lectureframe{00:13:30}
\label{fig:l03-test-function}
\end{figure}

One does not verify Eq.~\eqref{eq:l03-fundamental-lemma-hypothesis} by trying
functions one after another; that would take forever. In an application we
prove the integral vanishes for an arbitrary \(f\), and the proposition then
gives the pointwise conclusion in one step.

\section{Histories: Local and Global Descriptions}

Let a system have generalized coordinates
\begin{equation}
 q_1(t),q_2(t),\ldots,q_n(t).
\end{equation}
The complete collection of these functions is a \emph{history}, or a
trajectory in configuration space. For a system of \(N\) unconstrained
particles in three dimensions, \(n=3N\).

Whether time is drawn horizontally, as the independent variable often is in
elementary mathematics, or vertically, as it often is in physics diagrams, is
only a convention. The object represented is the same set of functions
\(q_i(t)\).

There are two apparently different ways to state a law of motion. A local law
uses information near one instant to continue the trajectory. In Newtonian
mechanics, knowledge of \(q_i\) and \(\dot q_i\), together with equations of
the form \(F=ma\), determines the acceleration and advances the solution to the
next instant.

A global law considers the whole path between two endpoint configurations. It
assigns a number, the action, to every candidate history and says that the
physical history makes that number stationary. Instead of giving \(q_i\) and
\(\dot q_i\) at one time, this boundary-value formulation fixes \(q_i\) at an
initial and a final time.

The two formulations must be equivalent when they describe the same physics.
There is already a simple reason to expect locality to emerge. If a path is
stationary on a long interval, its segment between any two neighboring points
must also be stationary with those points held fixed. Otherwise that segment
could be replaced and the action of the whole path would change to first
order.

\section{Varying a Path}

Write the physical history as \(\hat q_i(t)\). A nearby history can be
parameterized by
\begin{equation}
 q_i(t,\alpha)=\hat q_i(t)+\alpha f_i(t),
 \qquad
 f_i(t_1)=f_i(t_2)=0.
 \label{eq:l03-varied-path}
\end{equation}
The number \(\alpha\) controls the size and sign of the deformation; it does
not depend on time. The functions \(f_i(t)\) describe its shape and are
otherwise arbitrary. The endpoint conditions nail the beginning and end of
every candidate path in place.

\begin{figure}[H]
\centering
\includegraphics[width=0.79\textwidth]{lecture_03_figure_07.png}
\caption{A fixed-endpoint deformation at 00:22:55. The red curve is displaced
in the interior while the endpoint conditions \(f_i(t_1)=f_i(t_2)=0\) remain
visible on the board.}
\lectureframe{00:22:55}
\label{fig:l03-varied-path}
\end{figure}

For fixed \(\hat q_i\) and fixed deformation functions \(f_i\), the action is
now an ordinary function \(A(\alpha)\). The physical path corresponds to
\(\alpha=0\). If it is a minimum, then in particular it is stationary:
\begin{equation}
 \left.\frac{dA}{d\alpha}\right|_{\alpha=0}=0.
 \label{eq:l03-stationarity}
\end{equation}
The first derivative is the condition we need; whether the stationary path is
a minimum, maximum, or saddle requires additional information.

\begin{classroomqa}
\audiencequestion{What is the action being varied?}
\lecturerresponse{For a history between \(t_1\) and \(t_2\), the action is
\begin{equation}
 A[q]=\int_{t_1}^{t_2}L(q,\dot q,t)\,dt.
 \label{eq:l03-action}
\end{equation}
The Lagrangian \(L\) is evaluated at every point of the path. In the Newtonian
examples below it is \(L=T-U\), but the variational calculation needs only that
\(L\) be a function of the coordinates, velocities, and possibly time.}
\end{classroomqa}

\begin{figure}[H]
\centering
\includegraphics[width=0.82\textwidth]{lecture_03_figure_08.png}
\caption{The action integral at 00:28:42, written before any particular
mechanical form of the Lagrangian is needed.}
\lectureframe{00:28:42}
\label{fig:l03-action-board}
\end{figure}

Equation~\eqref{eq:l03-varied-path} immediately gives
\begin{equation}
 \frac{\partial q_i}{\partial\alpha}=f_i,
 \qquad
 \frac{\partial\dot q_i}{\partial\alpha}=\dot f_i.
\end{equation}
Applying the chain rule inside the action,
\begin{align}
 \frac{dA}{d\alpha}
 &=
 \int_{t_1}^{t_2}
 \sum_i
 \left(
   \frac{\partial L}{\partial q_i}
   \frac{\partial q_i}{\partial\alpha}
   +
   \frac{\partial L}{\partial\dot q_i}
   \frac{\partial\dot q_i}{\partial\alpha}
 \right)dt \\
 &=
 \int_{t_1}^{t_2}
 \sum_i
 \left(
   \frac{\partial L}{\partial q_i}f_i
   +
   \frac{\partial L}{\partial\dot q_i}\dot f_i
 \right)dt.
 \label{eq:l03-first-variation-before-parts}
\end{align}

\begin{figure}[H]
\centering
\includegraphics[width=0.86\textwidth]{lecture_03_figure_09.png}
\caption{The first variation at 00:34:10. Changes of both \(q_i\) and
\(\dot q_i\) contribute, and the audience supplies the required sum over
coordinates.}
\lectureframe{00:34:10}
\label{fig:l03-first-variation}
\end{figure}

\begin{classroomqa}
\audiencequestion{Should there be a sum over \(i\)?}
\lecturerresponse{Yes. Each coordinate and its velocity are independent
arguments of \(L\), so the chain rule contains one pair of terms for every
\(i\). For a one-coordinate problem the sum is invisible; for a general
system it is essential.}
\end{classroomqa}

The unwanted \(\dot f_i\) in the second term is precisely why integration by
parts was prepared. Applying Eq.~\eqref{eq:l03-integration-by-parts},
\begin{align}
 \int_{t_1}^{t_2}
 \frac{\partial L}{\partial\dot q_i}\dot f_i\,dt
 &=
 \left[
 f_i\frac{\partial L}{\partial\dot q_i}
 \right]_{t_1}^{t_2}
 -
 \int_{t_1}^{t_2}
 f_i\frac{d}{dt}
 \left(\frac{\partial L}{\partial\dot q_i}\right)dt.
\end{align}
In the earlier notation, \(f_i\) plays the role of \(f\), while
\(\partial L/\partial\dot q_i\) plays the role of \(g\). The operation removes
the dot from \(f_i\) and differentiates the momentum-like factor instead.
The boundary term vanishes because \(f_i(t_1)=f_i(t_2)=0\). Consequently,
\begin{equation}
 \left.\frac{dA}{d\alpha}\right|_{\alpha=0}
 =
 \int_{t_1}^{t_2}\sum_i
 \left[
   \frac{\partial L}{\partial q_i}
   -
   \frac{d}{dt}
   \left(\frac{\partial L}{\partial\dot q_i}\right)
 \right]_{\hat q}
 f_i(t)\,dt.
 \label{eq:l03-first-variation-after-parts}
\end{equation}

\begin{classroomqa}
\audiencequestion{What happened to the endpoint contribution?}
\lecturerresponse{It is exactly
\(\left[f_i\,\partial L/\partial\dot q_i\right]_{t_1}^{t_2}\).
The varied and physical paths have the same endpoints, so every \(f_i\)
vanishes there and the contribution is zero. The endpoint condition was chosen
at the start for this reason.}
\end{classroomqa}

Stationarity requires Eq.~\eqref{eq:l03-first-variation-after-parts} to vanish
for every collection of functions \(f_i(t)\). The arbitrary-function
proposition applies separately to every \(i\), giving
\begin{equation}
 \boxed{
 \frac{d}{dt}
 \left(\frac{\partial L}{\partial\dot q_i}\right)
 =
 \frac{\partial L}{\partial q_i}
 }
 \qquad (i=1,\ldots,n).
 \label{eq:l03-euler-lagrange}
\end{equation}
These are the Euler--Lagrange equations. The global condition
\(\delta A=0\) has become one local differential equation for each generalized
coordinate.

\begin{figure}[H]
\centering
\includegraphics[width=0.72\textwidth]{lecture_03_figure_02.png}
\caption{The Euler--Lagrange time-derivative term at 00:45:02, beside the
trajectory picture used to connect the global path with its local equation.}
\lectureframe{00:45:02}
\label{fig:l03-euler-lagrange-board}
\end{figure}

The result is not confined to point particles. Field theory, general
relativity, and string theory use the same variational structure with more
elaborate coordinates. Quantum mechanics also retains it in another form:
Feynman's path integral assigns amplitudes to histories through their action.

\section{Momentum and Generalized Force}

The quantity differentiated with respect to time in
Eq.~\eqref{eq:l03-euler-lagrange} has a name:
\begin{equation}
 \pi_i\equiv\frac{\partial L}{\partial\dot q_i}.
 \label{eq:l03-canonical-momentum}
\end{equation}
It is the canonical momentum conjugate to \(q_i\). The corresponding
\(\partial L/\partial q_i\) is the generalized force. The equation of motion
therefore reads
\begin{equation}
 \dot\pi_i=\frac{\partial L}{\partial q_i}.
 \label{eq:l03-momentum-force}
\end{equation}

\begin{figure}[H]
\centering
\includegraphics[width=0.64\textwidth]{lecture_03_figure_03.png}
\caption{Canonical momentum introduced on the board at 00:46:32. The standard
lowercase notation is
\(\pi_i=\partial L/\partial\dot q_i\).}
\lectureframe{00:46:32}
\label{fig:l03-canonical-momentum-board}
\end{figure}

\begin{classroomqa}
\audiencequestion{What does the Euler--Lagrange equation say in ordinary
language?}
\lecturerresponse{First, it says that a global statement about the entire
trajectory is equivalent to a local differential law. Second, after naming
the two sides, it says that the time derivative of momentum equals generalized
force. The concrete examples now show that this is \(F=ma\), written in a form
that survives arbitrary choices of coordinates.}
\end{classroomqa}

One may try to picture the left side as a time change of the
velocity-derivative of \(L\), and the right side as the coordinate-derivative
at the present point. That intuition is useful, but the names momentum and
force become convincing only after the first calculation.

\section{Recovering Newton's Equation}

For one particle of mass \(m\) moving on a line,
\begin{equation}
 L(x,\dot x)=\frac12m\dot x^2-U(x).
 \label{eq:l03-one-particle-lagrangian}
\end{equation}
The generalized coordinate \(q\) is simply \(x\), and its canonical momentum
is the ordinary momentum:
\begin{equation}
 \frac{\partial L}{\partial\dot x}=m\dot x=p.
\end{equation}
The coordinate derivative is
\begin{equation}
 \frac{\partial L}{\partial x}=-\frac{dU}{dx}.
\end{equation}
Hence
\begin{equation}
 m\ddot x=-\frac{dU}{dx},
\end{equation}
which is Newton's law for a conservative force.

The minus sign is not optional. Had we chosen \(L=T+U\), the force would have
come out with the wrong sign. The Lagrangian \(T-U\) is not the energy and is
not generally conserved along a trajectory. For a time-independent
conservative system the familiar conserved energy is \(T+U\); its derivation
from symmetry is deliberately postponed.

For many particles, combine every particle label and Cartesian direction into
one index \(i\). Then
\begin{equation}
 L=\sum_i\frac12m_i\dot x_i^2-U(x_1,\ldots,x_n),
\end{equation}
and, coordinate by coordinate,
\begin{equation}
 p_i=m_i\dot x_i,
 \qquad
 \dot p_i=-\frac{\partial U}{\partial x_i}.
 \label{eq:l03-many-particle-newton}
\end{equation}
For a physical particle the same mass occurs in its three Cartesian
components, although the notation is general enough to attach a coefficient
to every coordinate. A rigid body, an eraser, or a person can ultimately be
viewed as a large collection of particles held together by internal forces,
so this construction already covers an enormous class of mechanical systems.

\section{Translation Symmetry and Momentum}

Consider two interacting particles on a line:
\begin{equation}
 L=
 \frac12m_1\dot x_1^2+
 \frac12m_2\dot x_2^2-
 U(x_1-x_2).
 \label{eq:l03-two-particle-lagrangian}
\end{equation}
The potential depends only on their separation. If both coordinates are
shifted by the same constant \(a\),
\begin{equation}
 x_1\longmapsto x_1+a,
 \qquad
 x_2\longmapsto x_2+a,
\end{equation}
their velocities and the difference \(x_1-x_2\) are unchanged. The
Lagrangian, and therefore the action, is invariant. This is translation
symmetry.

\begin{figure}[htbp]
\centering
\begin{tikzpicture}[>=Stealth,thick]
  \draw (-3.4,0) -- (3.4,0);
  \fill (-1.55,0) circle (1.8pt)
    node[above=5pt] {$m_1$}
    node[below=5pt] {$x_1$};
  \fill (1.15,0) circle (1.8pt)
    node[above=5pt] {$m_2$}
    node[below=5pt] {$x_2$};
  \draw[<->] (-1.55,-0.75) -- (1.15,-0.75)
    node[midway,below] {$x_2-x_1$};
  \draw[->,blue!65!black] (-2.4,1.0) -- (-1.2,1.0)
    node[midway,above] {$a$};
  \draw[->,blue!65!black] (0.3,1.0) -- (1.5,1.0);
\end{tikzpicture}
\caption{A rigid translation moves both particles by the same amount without
changing their separation or velocities.}
\label{fig:l03-rigid-translation}
\end{figure}

This form would fail if a third, external lump of matter exerted a
position-dependent force. The potential would then depend on the distance from
that external object as well as on the separation between the two particles.
The two-particle subsystem would no longer be closed or translation invariant.

Set \(d=x_1-x_2\). The chain rule gives
\begin{align}
 \frac{\partial U}{\partial x_1}
 &=\frac{dU}{dd},
 &
 \frac{\partial U}{\partial x_2}
 &=-\frac{dU}{dd}.
 \label{eq:l03-opposite-potential-derivatives}
\end{align}
Therefore
\begin{align}
 \dot p_1&=-\frac{\partial U}{\partial x_1},
 &
 \dot p_2&=-\frac{\partial U}{\partial x_2}
          =+\frac{\partial U}{\partial x_1}.
\end{align}
The forces are equal and opposite, and adding the equations gives
\begin{equation}
 \frac{d}{dt}(p_1+p_2)=0.
 \label{eq:l03-total-momentum}
\end{equation}

\begin{figure}[H]
\centering
\includegraphics[width=0.86\textwidth]{lecture_03_figure_10.png}
\caption{The translation-invariant Lagrangian and its two force equations at
01:08:50. Dependence on \(x_1-x_2\) makes the derivatives equal and opposite.}
\lectureframe{01:08:50}
\label{fig:l03-translation-board}
\end{figure}

This is the first instance of the central pattern:
\[
 \text{translation symmetry}
 \quad\Longrightarrow\quad
 \text{conservation of linear momentum}.
\]
The symmetry is an active operation on the system that leaves the action
unchanged. Here the operation can be made by an arbitrarily small amount and
built up continuously, so it is a continuous symmetry.

\begin{classroomqa}
\audiencequestion{If \(U\) depends on \(\lvert x_1-x_2\rvert\), is inversion
another symmetry?}
\lecturerresponse{Yes. Interchanging the orientation of the separation leaves
such a potential unchanged. Unlike translation, inversion is discrete: there
is no continuously adjustable parameter connecting the two alternatives.
Continuous symmetries generate the classical conservation laws being studied
here. Discrete symmetries become especially consequential in quantum
mechanics, which we leave for later.}
\end{classroomqa}

\section{One Translation Survives}

Place a particle near Earth's surface and retain horizontal \(x\) and vertical
\(y\) motion. With \(y\) measured upward,
\begin{equation}
 L=\frac12m\dot x^2+\frac12m\dot y^2-mgy,
 \label{eq:l03-near-earth-lagrangian}
\end{equation}
where \(g\) is the approximately constant local gravitational acceleration.
Numerically it is about \(9.8\,\mathrm{m\,s^{-2}}\), commonly rounded to
\(10\,\mathrm{m\,s^{-2}}\) for this discussion.

\begin{figure}[H]
\centering
\includegraphics[width=0.78\textwidth]{lecture_03_figure_04.png}
\caption{The two-coordinate kinetic energy and near-Earth setup at 01:15:15.
The stable board content is \(T=\tfrac12m\dot x^2+\tfrac12m\dot y^2\);
the potential is completed immediately afterward.}
\lectureframe{01:15:15}
\label{fig:l03-near-earth-setup}
\end{figure}

The Lagrangian contains no \(x\), so horizontal translation is a symmetry:
\begin{equation}
 \frac{d}{dt}(m\dot x)=0.
\end{equation}
It does contain \(y\), so vertical translation is not a symmetry:
\begin{equation}
 \frac{d}{dt}(m\dot y)=-mg,
 \qquad
 \ddot y=-g.
 \label{eq:l03-near-earth-equations}
\end{equation}
There is one translation symmetry and one corresponding conserved momentum,
not two.

\begin{figure}[H]
\centering
\includegraphics[width=0.82\textwidth]{lecture_03_figure_11.png}
\caption{The completed near-Earth Lagrangian and coordinate derivatives at
01:19:15. The absence of \(x\) gives horizontal momentum conservation, while
\(\partial U/\partial y=mg\) gives downward acceleration.}
\lectureframe{01:19:15}
\label{fig:l03-near-earth-equations}
\end{figure}

\begin{classroomqa}
\audiencequestion{For a very small vertical displacement, is there an
approximate \(y\)-translation symmetry?}
\lecturerresponse{No exact symmetry appears merely because the displacement is
small. In the near-Earth approximation \(U=mgy\), any nonzero change of \(y\)
changes the potential by \(mg\,\Delta y\). We may add one fixed constant to
\(U\) and thereby choose a different zero of potential energy, but that does
not remove the \(y\)-dependence or the vertical force.}
\end{classroomqa}

Adding \(C\) to the potential changes neither force nor trajectory:
\begin{equation}
 U(y)\longmapsto U(y)+C,
 \qquad
 -\frac{d}{dy}\bigl(U+C\bigr)=-\frac{dU}{dy}.
\end{equation}
Sea level, the floor, or any other fixed height may therefore be chosen as the
zero of potential energy.

\begin{classroomqa}
\audiencequestion{Is this just the familiar falling-body problem and its
energy conservation in another form?}
\lecturerresponse{It is the same mechanics, but so far we have derived the
equations of motion and identified momentum conservation. The conserved
energy is \(T+U\), not the Lagrangian \(T-U\). Its origin will be traced to
another symmetry, time translation, in a later step.}
\end{classroomqa}

\section{Changing Coordinates}

The action does not depend on how we label the configurations. Rectangular,
polar, parabolic, or curved coordinates may make the formulas look different,
but a correctly transformed action assigns the same number to the same
physical history. Stationarity of that action is therefore
coordinate-independent.

For planar motion, replace \(x,y\) by polar coordinates \(r,\theta\). At every
instant the velocity has a radial component and a perpendicular component:
\begin{equation}
 v_r=\dot r,
 \qquad
 v_\perp=r\dot\theta.
 \label{eq:l03-polar-velocity}
\end{equation}
The factor of \(r\) has a simple geometric origin. A fixed angular change
\(d\theta\) sweeps the arc length \(r\,d\theta\), so the same angular speed
corresponds to a larger linear speed farther from the origin. The fingertips
of a rotating arm move faster than the elbow.

\begin{figure}[H]
\centering
\includegraphics[width=0.83\textwidth]{lecture_03_figure_12.png}
\caption{Polar-coordinate velocity components at 01:27:50. The board records
\(v_r=\dot r\), \(v_\perp=r\dot\theta\), and the resulting kinetic-energy
terms.}
\lectureframe{01:27:50}
\label{fig:l03-polar-kinetic-board}
\end{figure}

Hence
\begin{equation}
 T=\frac12m\left(\dot r^2+r^2\dot\theta^2\right).
 \label{eq:l03-polar-kinetic}
\end{equation}
For a central potential \(U(r)\),
\begin{equation}
 L=
 \frac12m\dot r^2+
 \frac12mr^2\dot\theta^2-
 U(r).
 \label{eq:l03-central-lagrangian}
\end{equation}
It is useful to write both Euler--Lagrange equations. The radial one is
\begin{equation}
 m\ddot r
 =
 mr\dot\theta^2-\frac{dU}{dr}.
 \label{eq:l03-radial-equation}
\end{equation}
The \(mr\dot\theta^2\) term comes from differentiating the angular kinetic
energy with respect to \(r\). Radial momentum is not conserved: the central
force and the changing radial direction both enter its evolution.

The angular coordinate is different. The Lagrangian depends on
\(\dot\theta\) but not on \(\theta\) itself. Thus \(\theta\) is a cyclic
coordinate and
\begin{equation}
 p_\theta
 =
 \frac{\partial L}{\partial\dot\theta}
 =
 mr^2\dot\theta.
 \label{eq:l03-angular-momentum}
\end{equation}
Its equation of motion is
\begin{equation}
 \frac{d}{dt}\left(mr^2\dot\theta\right)
 =
 \frac{\partial L}{\partial\theta}
 =
 0.
 \label{eq:l03-angular-conservation}
\end{equation}

\begin{figure}[H]
\centering
\includegraphics[width=0.60\textwidth]{lecture_03_figure_13.png}
\caption{Angular momentum conservation completed on the board at 01:33:35:
\(\frac{d}{dt}(mr^2\dot\theta)=0\).}
\lectureframe{01:33:35}
\label{fig:l03-angular-momentum-board}
\end{figure}

The conserved quantity \(mr^2\dot\theta\) is the angular momentum, usually
denoted by \(\ell\) when \(L\) is already reserved for the Lagrangian. Solving
for the angular speed,
\begin{equation}
 \dot\theta=\frac{\ell}{mr^2}.
\end{equation}
As \(r\) decreases, \(\dot\theta\) increases. This is the spinning-skater
effect in its simplest form.

We have now seen the same structure twice:
\begin{align*}
 \text{translation invariance}
 &\quad\Longrightarrow\quad
 \text{linear momentum conservation},\\
 \text{rotational invariance}
 &\quad\Longrightarrow\quad
 \text{angular momentum conservation}.
\end{align*}
The action principle makes the pattern visible. A coordinate that does not
appear in the Lagrangian has a conserved conjugate momentum, and the local
equation expressing that fact is itself the consequence of stationarity of
the entire history.
